\subsection{Use Case Diagram}
\label{sec:usecases-revised}

แผนภาพต่อไปนี้แบ่งตามกลุ่มงานของระบบ Pacegasus โดยใช้รหัส UC-01 ถึง UC-20 ตรงกับตารางคำอธิบายทุกกรณี นักวิ่งหมายถึงผู้ใช้แอปพลิเคชันทั้งก่อนและหลังยืนยันตัวตน โดยแต่ละกรณีระบุเงื่อนไขเข้าสู่ระบบแยกกัน หัวหน้าชมรมเป็นบทบาทเฉพาะของนักวิ่ง จึงรับความสามารถของนักวิ่งผ่าน Actor Generalization และมีสิทธิ์จัดการสมาชิกเฉพาะชมรมของตนเอง สมาชิกชมรมใช้ Actor นักวิ่งร่วมกับเงื่อนไขสมาชิกภาพ ไม่ได้หมายความว่านักวิ่งทุกคนเป็นสมาชิกทุกชมรม

แผนภาพและตารางเป็นข้อกำหนดการออกแบบ ไม่ใช่การรับรองว่าทุกเส้นทางพัฒนาและทดสอบครบแล้ว UC-20 เป็นส่วนต่อยอดสำหรับ Strava ส่วน UC-07, UC-14--16, UC-18--19 ยังต้องตรวจสอบรายละเอียดหรือทดสอบเพิ่มเติมตามข้อจำกัดที่ระบุในตาราง ผลทดสอบจริงแสดงแยกในหัวข้อการทดสอบระบบ ไม่มีการใช้ความสัมพันธ์ include หรือ extend แทนลำดับหน้าจอ; การเรียกใช้กรณีอื่นระบุไว้ในขั้นตอนที่เกี่ยวข้อง

\newcommand{\UCActor}[3]{%
\draw (#1,#2+0.3) circle (0.12);
\draw (#1,#2+0.18)--(#1,#2-0.22);
\draw (#1-0.22,#2)--(#1+0.22,#2);
\draw (#1,#2-0.22)--(#1-0.2,#2-0.5);
\draw (#1,#2-0.22)--(#1+0.2,#2-0.5);
\node[font=\small,align=center] at (#1,#2-0.8) {#3};}


\begin{figure}[H]\centering

\begin{tikzpicture}[x=1cm,y=1cm,uc/.style={ellipse,draw,fill=green!8,text width=4.5cm,minimum height=1.05cm,align=center,font=\small}]

\draw (1,0) rectangle (8.3,-8.8);

\node[font=\small\bfseries,text width=7cm,align=center] at (4.65,-0.5) {Pacegasus: บัญชีและโปรไฟล์};

\UCActor{0}{-4.4}{นักวิ่ง}

\node[uc] (u1) at (4.65,-1.8) {UC-01 จัดการการเข้าใช้งานระบบ};
\draw (0.25,-4.4) -- (u1.west);
\node[uc] (u5) at (4.65,-4.0) {UC-02 จัดทำโปรไฟล์ครั้งแรก};
\draw (0.25,-4.4) -- (u5.west);
\node[uc] (u6) at (4.65,-6.2) {UC-03 แก้ไขโปรไฟล์};
\draw (0.25,-4.4) -- (u6.west);

\UCActor{9.5}{-1.4}{บริการยืนยัน\\อีเมล}
\draw (9.25,-1.4)--(u1.east);
\UCActor{9.5}{-4.1}{Google}
\draw (9.25,-4.1)--(u1.east);

\end{tikzpicture}

\caption{Use Case Diagram: บัญชีและโปรไฟล์}\label{fig:usecase-group-1}

\end{figure}

\begin{figure}[H]\centering

\begin{tikzpicture}[x=1cm,y=1cm,uc/.style={ellipse,draw,fill=green!8,text width=4.5cm,minimum height=1.05cm,align=center,font=\small}]

\draw (1,0) rectangle (8.3,-8.8);

\node[font=\small\bfseries,text width=7cm,align=center] at (4.65,-0.5) {Pacegasus: สุขภาวะและแผนฝึก};

\UCActor{0}{-4.4}{นักวิ่ง}

\node[uc] (u7) at (4.65,-1.8) {UC-04 บันทึกสุขภาวะประจำวัน};

\draw (0.25,-4.4) -- (u7.west);

\node[uc] (u8) at (4.65,-3.8) {UC-05 เลือกแผนการฝึก};

\draw (0.25,-4.4) -- (u8.west);

\node[uc] (u9) at (4.65,-5.8) {UC-06 จัดตารางฝึก};

\draw (0.25,-4.4) -- (u9.west);

\node[uc] (u10) at (4.65,-7.8) {UC-07 ดูคำแนะนำและผลปรับแผน};

\draw (0.25,-4.4) -- (u10.west);

\end{tikzpicture}

\caption{Use Case Diagram: สุขภาวะและแผนฝึก}\label{fig:usecase-group-2}

\end{figure}

\begin{figure}[H]\centering

\begin{tikzpicture}[x=1cm,y=1cm,uc/.style={ellipse,draw,fill=green!8,text width=4.5cm,minimum height=1.05cm,align=center,font=\small}]

\draw (1,0) rectangle (8.3,-10.8);

\node[font=\small\bfseries,text width=7cm,align=center] at (4.65,-0.5) {Pacegasus: การวิ่งและภารกิจ};

\UCActor{0}{-5.4}{นักวิ่ง}

\node[uc] (u11) at (4.65,-1.8) {UC-08 ดำเนินกิจกรรมการวิ่ง};
\draw (0.25,-5.4) -- (u11.west);
\node[uc] (u14) at (4.65,-3.8) {UC-09 ดูผลลัพธ์และประวัติการวิ่ง};
\draw (0.25,-5.4) -- (u14.west);
\node[uc] (u15) at (4.65,-5.8) {UC-10 บันทึกการประเมินหลังวิ่ง};
\draw (0.25,-5.4) -- (u15.west);
\node[uc] (u16) at (4.65,-7.8) {UC-11 เลือกบทและเริ่มภารกิจหลัก};
\draw (0.25,-5.4) -- (u16.west);
\node[uc] (u17) at (4.65,-9.8) {UC-12 สิ้นสุดภารกิจและรับรางวัล};
\draw (0.25,-5.4) -- (u17.west);

\end{tikzpicture}

\caption{Use Case Diagram: การวิ่งและภารกิจ}\label{fig:usecase-group-3}

\end{figure}

\begin{figure}[H]\centering

\begin{tikzpicture}[x=1cm,y=1cm,uc/.style={ellipse,draw,fill=green!8,text width=4.5cm,minimum height=1.05cm,align=center,font=\small}]

\draw (1,0) rectangle (8.3,-8.8);

\node[font=\small\bfseries,text width=7cm,align=center] at (4.65,-0.5) {Pacegasus: สังคมและชมรม};

\UCActor{0}{-3.1}{นักวิ่ง}

\node[uc] (u18) at (4.65,-1.8) {UC-13 จัดการชมรม};

\draw (0.25,-3.1) -- (u18.west);

\node[uc] (u19) at (4.65,-3.8) {UC-14 เข้าร่วมหรือออกจากชมรม};

\draw (0.25,-3.1) -- (u19.west);

\node[uc] (u20) at (4.65,-5.8) {UC-15 จัดการสมาชิกในชมรม};

\node[uc] (u21) at (4.65,-7.8) {UC-16 จัดการรายชื่อเพื่อน};

\draw (0.25,-3.1) -- (u21.west);

\UCActor{0}{-7.0}{หัวหน้าชมรม}

\draw[-{Triangle[open,length=3mm,width=3mm]}] (-0.22,-7.0)--(-0.55,-7.0)--(-0.55,-3.1)--(-0.22,-3.1);

\draw (0.25,-7.0)--(u20.west);

\end{tikzpicture}

\caption{Use Case Diagram: สังคมและชมรม}\label{fig:usecase-group-4}

\end{figure}

\begin{figure}[H]\centering

\begin{tikzpicture}[x=1cm,y=1cm,uc/.style={ellipse,draw,fill=green!8,text width=4.5cm,minimum height=1.05cm,align=center,font=\small}]

\draw (1,0) rectangle (8.3,-8.8);

\node[font=\small\bfseries,text width=7cm,align=center] at (4.65,-0.5) {Pacegasus: ไอเทม อวตาร์ อันดับ และบริการภายนอก};

\UCActor{0}{-4.4}{นักวิ่ง}

\node[uc] (u22) at (4.65,-1.8) {UC-17 ซื้อไอเทม};

\draw (0.25,-4.4) -- (u22.west);

\node[uc] (u23) at (4.65,-3.8) {UC-18 จัดการอวตาร์};

\draw (0.25,-4.4) -- (u23.west);

\node[uc] (u24) at (4.65,-5.8) {UC-19 ดูตารางจัดอันดับ};

\draw (0.25,-4.4) -- (u24.west);

\node[uc] (u25) at (4.65,-7.8) {UC-20 เชื่อมต่อและนำเข้าการวิ่งจาก Strava};

\draw (0.25,-4.4) -- (u25.west);

\UCActor{9.5}{-5.95}{Strava}

\draw (9.25,-5.95)--(u25.east);

\end{tikzpicture}

\caption{Use Case Diagram: ไอเทม อวตาร์ อันดับ และบริการภายนอก}\label{fig:usecase-group-5}

\end{figure}

\subsection{Use Case Description}
แต่ละตารางระบุรหัส เป้าหมาย Actor เงื่อนไขก่อนเริ่ม ผลสำเร็จ ขั้นตอนหลัก และทางเลือก/ข้อผิดพลาด กรณีที่มีบริการภายนอก นักวิ่งเป็น Primary Actor และบริการนั้นเป็น Supporting Actor ข้อกำหนดร่วมคือระบบต้องตรวจสิทธิ์ฝั่งระบบก่อนอ่านหรือแก้ข้อมูล ไม่ถือว่าการซ่อนปุ่มเป็นการตรวจสิทธิ์ และต้องไม่แจ้งสำเร็จเมื่อการบันทึกยังล้มเหลว

รางวัลเช็กอิน รางวัลประเมินหลังวิ่ง และรางวัลภารกิจเป็นคนละเหตุการณ์ ต้องอ้างอิงผู้ใช้ วันที่หรือรหัสเซสชัน/ภารกิจ และประเภทรางวัลเพื่อป้องกันการให้ซ้ำ การสิ้นสุดเซสชันไม่ใช่เหตุให้แจกซ้ำกับรางวัลภารกิจ หน่วยกิตนับเฉพาะผลที่ผ่านเกณฑ์แผนและแยกตามแผน เมื่อครบเกณฑ์จึงปลดล็อกระดับถัดไป ข้อกำหนดการป้องกันซ้ำเป็นเงื่อนไขที่ต้องตรวจสอบ ไม่ใช่ข้ออ้างว่าทุก endpoint ผ่านการทดสอบแล้ว

\newcommand{\UseCaseTable}[6]{%
\clearpage
\begin{longtable}{|>{\raggedright\arraybackslash}p{0.23\textwidth}|>{\raggedright\arraybackslash}p{0.69\textwidth}|}
\caption{Use Case Description: #1}\\\hline
\textbf{Use Case} & #1 \\\hline
\endfirsthead
\multicolumn{2}{l}{\small\textbf{ตารางที่~\thetable} Use Case Description: #1 (ต่อ)}\\\noalign{\vskip12pt}\hline
\endhead
\textbf{Actor} & #2 \\\hline
\textbf{Precondition} & #3 \\\hline
\textbf{Postcondition} & #4 \\\hline
\textbf{Main Flow} & #5 \\\hline
\textbf{Alternative / Exception} & #6 \\\hline
\end{longtable}}


\UseCaseTable{UC-01 จัดการการเข้าใช้งานระบบ}
{นักวิ่ง; บริการยืนยันอีเมล; Google}
{ผู้ใช้ต้องการสมัครสมาชิก เข้าสู่ระบบ หรือออกจากระบบ}
{บัญชีถูกสร้างหรือยืนยันตัวตนและเข้าสู่ระบบสำเร็จ หรือเซสชันถูกยกเลิกเมื่อออกจากระบบ}
{1) ผู้ใช้เลือกสมัครสมาชิก เข้าสู่ระบบด้วย Email OTP หรือ Google หรือออกจากระบบ\newline 2) กรณี Email OTP ระบบตรวจอีเมลและให้บริการยืนยันอีเมลส่งรหัส; ผู้ใช้กรอกรหัสเพื่อยืนยัน\newline 3) กรณี Google ผู้ใช้เลือกบัญชีและระบบตรวจผลยืนยันจาก Google\newline 4) ระบบสร้างบัญชีใหม่เมื่อจำเป็น ตรวจสถานะโปรไฟล์ แล้วเข้าสู่หน้าหลักหรือทำ UC-02 ต่อ\newline 5) กรณีออกจากระบบ ระบบขอคำยืนยัน ยกเลิกเซสชัน และกลับสู่หน้าเข้าสู่ระบบ}
{A/E1: ไม่ยอมรับข้อกำหนดที่จำเป็น ไม่สร้างบัญชี\newline A/E2: OTP ผิด หมดอายุ หรือส่งไม่สำเร็จ ให้ลองใหม่ตามข้อจำกัดโดยไม่ถือว่ายืนยันสำเร็จ\newline A/E3: Google ยกเลิกหรือยืนยันไม่สำเร็จ ไม่สร้างสถานะเข้าสู่ระบบ; บัญชีเดิมต้องยืนยันความเป็นเจ้าของก่อนผูก Google\newline A/E4: ผู้ใช้ยกเลิกการออกจากระบบ ให้คงเซสชันเดิม; ยกเลิก token ไม่สำเร็จให้แจ้งข้อผิดพลาด}


\UseCaseTable{UC-02 จัดทำโปรไฟล์ครั้งแรก}{นักวิ่ง}{เข้าสู่ระบบแล้วและ Onboarding ยังไม่สมบูรณ์}{ข้อมูลพื้นฐาน สุขภาพ ประวัติวิ่ง และเป้าหมายครบถ้วน พร้อมสถานะ Onboarding สำเร็จ}{1) ระบบอ่านขั้นล่าสุดที่บันทึกไว้\newline 2) ผู้ใช้กรอกวันเกิด เพศ น้ำหนัก ส่วนสูง และจำนวนวันที่ต้องการวิ่ง\newline 3) ผู้ใช้ระบุโรคประจำตัวและอาการบาดเจ็บ หรือเลือกไม่มี\newline 4) ผู้ใช้กรอกประวัติวิ่งและเป้าหมายด้านระยะทางหรือสุขภาพ\newline 5) ระบบตรวจและบันทึกแต่ละขั้นตามลำดับ โดยยังไม่ถือว่าโปรไฟล์สมบูรณ์จนข้อมูลที่จำเป็นครบ\newline 6) ระบบกำหนดสถานะสมบูรณ์และแสดงหน้าหลัก}{A/E1: ขั้น 2--4: ข้อมูลไม่ครบหรือขัดกัน เช่น เลือกไม่มีพร้อมอาการบาดเจ็บ ให้แก้ช่องที่ผิด\newline A/E2: ขั้น 5: ออกจากขั้นตอนกลางคัน ให้กลับมาทำต่อจากขั้นที่บันทึกสำเร็จ; บันทึกล้มเหลวไม่เลื่อนสถานะขั้น\newline A/E3: ผู้ไม่มีประวัติวิ่งเลือกไม่มีประสบการณ์ได้ ไม่บังคับให้สร้างประวัติขึ้นเอง}


\UseCaseTable{UC-03 แก้ไขโปรไฟล์}{นักวิ่ง}{เข้าสู่ระบบและมีโปรไฟล์แล้ว}{ข้อมูลที่ยืนยันถูกบันทึกในโปรไฟล์ของตนเอง}{1) เปิดโปรไฟล์และเลือกแก้ไข\newline 2) ระบบแสดงข้อมูลปัจจุบันที่อนุญาตให้แก้\newline 3) ผู้ใช้แก้ข้อมูลพื้นฐาน สุขภาพ หรือเป้าหมาย แล้วบันทึก\newline 4) ระบบตรวจสอบและบันทึกข้อมูล\newline 5) ระบบแสดงค่าล่าสุด และใช้ข้อมูลใหม่ในการประเมินแผนครั้งถัดไปตาม UC-07}{A/E1: ขั้น 3: ยกเลิกให้คงค่าเดิม\newline A/E2: ขั้น 4: ข้อมูลไม่ถูกต้องให้แก้ไข; บันทึกล้มเหลวให้ลองใหม่และไม่แสดงผลสำเร็จ\newline A/E3: การเปลี่ยนเป้าหมายไม่โอนหน่วยกิตหรือเปลี่ยนแผนที่ลงทะเบียนโดยอัตโนมัติ ต้องผ่าน UC-05}


\UseCaseTable{UC-04 บันทึกสุขภาวะประจำวัน}{นักวิ่ง}{เข้าสู่ระบบและยังไม่มี Check-in สำหรับวันปัจจุบันตามเขตเวลาที่ระบบใช้}{บันทึก Wellness ของวันนั้นและรางวัลเช็กอินเพียงครั้งเดียว}{1) เปิด Daily Wellness Check-in\newline 2) ระบบแสดงหัวข้อคุณภาพการนอน ระดับพลังงาน ความเมื่อยล้าของกล้ามเนื้อ ระดับความเครียด และแรงจูงใจ\newline 3) ผู้ใช้ให้คะแนนแต่ละหัวข้อในช่วง 1--5 และอาจเพิ่มหมายเหตุ\newline 4) ระบบตรวจข้อมูล คำนวณคะแนนสุขภาวะรวม และบันทึก Check-in\newline 5) ระบบบันทึกรางวัลเช็กอิน และนำข้อมูลไปใช้ใน UC-07}{A/E1: ขั้น 4: คะแนนนอกช่วง 1--5 หรือไม่ได้ให้คะแนนครบทั้ง 5 หัวข้อ ให้แก้ไข\newline A/E2: ขั้น 4: มีข้อมูลของวันนั้นแล้ว ให้แสดงรายการเดิมและไม่ให้รางวัลซ้ำ\newline A/E3: ขั้น 4--5: บันทึกล้มเหลวให้ลองใหม่ โดยตรวจรายการเดิมก่อนสร้างรางวัล}


\UseCaseTable{UC-05 เลือกแผนการฝึก}{นักวิ่ง}{เข้าสู่ระบบและ Onboarding สมบูรณ์}{มีแผนที่ลงทะเบียนและบัญชีหน่วยกิตแยกตามแผน}{1) เปิดรายการแผนการฝึก\newline 2) ระบบแสดงระดับที่ปลดล็อก เป้าหมาย และข้อกำหนดแต่ละแผน\newline 3) ผู้ใช้เลือกแผนและยืนยัน\newline 4) ระบบตรวจสิทธิ์ระดับแล้วสร้างหรือเปิดแผนที่ลงทะเบียนไว้\newline 5) ระบบแสดงความคืบหน้าและให้จัดตารางด้วย UC-06}{A/E1: ขั้น 3--4: ระดับยังล็อก ให้แสดงเงื่อนไขปลดล็อกและไม่ลงทะเบียน\newline A/E2: ขั้น 4: ลงทะเบียนแผนเดิมแล้วให้เปิดรายการเดิม ไม่สร้างซ้ำ\newline A/E3: หน่วยกิตไม่โอนข้ามแผนต่างระดับ; การเลือกเป้าหมาย Full Marathon ไม่ถือว่ามีแผน Full Marathon ซึ่งยังไม่มีรายละเอียดในขอบเขตนี้}


\UseCaseTable{UC-06 จัดตารางฝึก}{นักวิ่ง}{เข้าสู่ระบบและมีแผนการฝึกที่ลงทะเบียน}{บันทึกตารางที่ผ่านกฎระดับและช่วงของแผน}{1) เปิดตารางของแผนและสัปดาห์ที่ต้องการ\newline 2) ระบบแสดงวันและประเภทการฝึกที่เลือกได้\newline 3) ผู้ใช้เลือกหรือแก้วันและประเภทการฝึก\newline 4) ระบบตรวจ Weekly Cap วันพัก ลำดับ Easy ก่อน Long Run สำหรับ Beginner และข้อจำกัดช่วง Taper/Race ตามกฎแผน\newline 5) ผู้ใช้ยืนยันและระบบบันทึกตาราง}{A/E1: ขั้น 4: ผิดกฎ ให้ระบุเหตุผลและกลับไปเลือกวันหรือประเภทใหม่ โดยไม่บันทึกตารางที่ผิด\newline A/E2: ขั้น 5: ยกเลิกให้คงตารางเดิม; บันทึกล้มเหลวให้ลองใหม่\newline A/E3: ไม่แก้ผลการวิ่งที่สิ้นสุดแล้วผ่านการแก้ตาราง}


\UseCaseTable{UC-07 ดูคำแนะนำและผลปรับแผน}{นักวิ่ง}{เข้าสู่ระบบและมีแผนการฝึก}{ผู้ใช้เห็นคำแนะนำล่าสุดพร้อมเหตุผลและสถานะว่าตารางเปลี่ยนหรือไม่}{1) ผู้ใช้เปิดแผนรายวัน หลังบันทึก Wellness หรือข้อมูลหลังวิ่ง\newline 2) ระบบรวบรวมข้อมูลที่มีของผู้ใช้และประเมินด้วยกฎการฝึกที่กำหนด\newline 3) ระบบแสดงคำแนะนำฝึกตามแผน ลดภาระ หรือพัก พร้อมเหตุผล\newline 4) หากมีการปรับตาราง ระบบตรวจข้อจำกัดของ UC-06 ก่อนบันทึกและแสดงรายการที่เปลี่ยน}{A/E1: ขั้น 2: ข้อมูลยังไม่พอ ให้แสดงข้อจำกัดของการประเมินและไม่แสดงดัชนีที่คำนวณไม่ได้\newline A/E2: ขั้น 2--4: ประเมินหรือบันทึกไม่สำเร็จ ให้คงแผนเดิมและแจ้งว่าไม่มีการปรับสำเร็จ\newline A/E3: รายละเอียดอัลกอริทึมและการใช้ Machine Learning เป็นส่วนที่ต้องตรวจสอบเพิ่มเติม ไม่ถือว่าผลทดสอบ API ยืนยันการปรับแผนครบถ้วน}


\UseCaseTable{UC-08 ดำเนินกิจกรรมการวิ่ง}
{นักวิ่ง}
{ผู้ใช้เข้าสู่ระบบและเลือกประเภทการวิ่งแล้ว}
{เซสชันการวิ่งและความคืบหน้าภารกิจที่เกี่ยวข้องถูกบันทึก โดยสิ้นสุดเซสชันเมื่อผู้ใช้ยืนยัน}
{1) ผู้ใช้เริ่มวิ่ง ระบบตรวจความพร้อมและสิทธิ์ GPS เมื่อจำเป็น แล้วสร้างเซสชัน\newline 2) ระบบติดตามและแสดงเวลา ระยะทาง Pace และสถานะเซสชัน\newline 3) ผู้ใช้เลือกทำภารกิจระหว่างวิ่ง ระบบตรวจหลักฐานและบันทึกความคืบหน้า\newline 4) ผู้ใช้เลือกสิ้นสุดการวิ่งและยืนยัน ระบบหยุดติดตาม บันทึกข้อมูล และสรุปภารกิจที่ผูกไว้\newline 5) ระบบแสดงผลด้วย UC-09 และเปิดการประเมินหลังวิ่งด้วย UC-10}
{A/E1: ไม่พร้อมใช้งาน GPS หรือสร้างเซสชันไม่สำเร็จ ให้แจ้งและไม่เริ่มการติดตามกลางแจ้ง\newline A/E2: หลักฐานภารกิจไม่ครบ ให้แจ้งสิ่งที่ขาดและคงความคืบหน้าที่ตรวจสอบแล้ว\newline A/E3: ผู้ใช้ยกเลิกการจบวิ่ง ให้กลับไปยังเซสชันเดิม\newline A/E4: GPS ขาดช่วงหรือบันทึกล้มเหลว ให้ระบุสถานะข้อมูล ไม่สร้างระยะทางที่ไม่มีหลักฐาน และไม่ให้รางวัลซ้ำ}


\UseCaseTable{UC-09 ดูผลลัพธ์และประวัติการวิ่ง}{นักวิ่ง}{เข้าสู่ระบบและมีเซสชันของตนเองที่บันทึกแล้ว}{แสดงผลการวิ่งและสถานะความครบถ้วนของข้อมูล}{1) เปิดผลหลังจบวิ่งหรือเลือกจากประวัติ\newline 2) ระบบตรวจสิทธิ์และอ่านข้อมูลเซสชัน\newline 3) ระบบแสดงเวลา ระยะทาง Pace และสถิติอื่นที่มีข้อมูลเพียงพอ\newline 4) ผู้ใช้ดูรายละเอียดและสถานะภารกิจหรือหน่วยกิตที่เกี่ยวข้อง}{A/E1: ขั้น 2: ไม่พบข้อมูลหรือไม่มีสิทธิ์ ให้แจ้งและกลับหน้าประวัติ\newline A/E2: ขั้น 2: โหลดไม่สำเร็จให้ลองใหม่\newline A/E3: ขั้น 3: ข้อมูลไม่ครบให้แสดงข้อจำกัดและไม่แทนค่าที่ขาดด้วยผลคำนวณที่ไม่มีหลักฐาน}


\UseCaseTable{UC-10 บันทึกการประเมินหลังวิ่ง}{นักวิ่ง}{เข้าสู่ระบบและเซสชันของตนเองสิ้นสุดแล้ว}{บันทึก RPE และข้อมูลหลังวิ่งผูกกับเซสชัน พร้อมรางวัลการประเมินตามกฎ}{1) ระบบเปิดแบบประเมินของเซสชันที่สิ้นสุด\newline 2) ผู้ใช้ระบุระยะเวลาวิ่ง (นาที) คะแนน RPE 1--10 ระดับความเครียด 1--10 ความรู้สึกหลังวิ่ง (หมดแรง แย่ เฉยๆ ดี ดีมาก) และระบุว่ามีอาการเจ็บหรือไม่ พร้อมรายละเอียดเมื่อมี\newline 3) ผู้ใช้ยืนยัน\newline 4) ระบบตรวจ คำนวณค่า session RPE จากระยะเวลาคูณคะแนน RPE และบันทึกข้อมูลผูกกับเซสชัน\newline 5) ระบบบันทึกรางวัลการประเมินเพียงครั้งเดียว และใช้ข้อมูลประกอบ UC-07}{A/E1: ขั้น 4: คะแนนนอกช่วง ระยะเวลาไม่มากกว่า 0 หรือข้อมูลจำเป็นไม่ครบ ให้แก้ไข\newline A/E1.1: ขั้น 4: ระบุว่ามีอาการเจ็บแต่ไม่กรอกรายละเอียด ให้แจ้งให้กรอก และไม่เก็บรายละเอียดอาการเมื่อระบุว่าไม่มี\newline A/E2: ขั้น 4: เซสชันยังไม่จบหรือไม่ใช่ของผู้ใช้ ให้ปฏิเสธ\newline A/E3: ขั้น 4--5: ส่งซ้ำให้คืนผลรายการเดิม ไม่ให้รางวัลเดิมซ้ำ; บันทึกล้มเหลวให้ลองใหม่}


\UseCaseTable{UC-11 เลือกบทและเริ่มภารกิจหลัก}{นักวิ่ง}{เข้าสู่ระบบและมีสิทธิ์เข้าถึงบทตามระดับ}{เริ่มภารกิจหลักที่เลือกและติดตามความคืบหน้าได้}{1) เปิดรายการบทและสถานะปลดล็อก\newline 2) ผู้ใช้เลือกบทและอ่านภารกิจ\newline 3) ระบบตรวจระดับและเงื่อนไขเริ่ม\newline 4) ผู้ใช้เริ่มภารกิจ; หากต้องวิ่งให้เริ่มหรือผูกเซสชันตาม UC-08\newline 5) ระบบบันทึกความคืบหน้าตามเกณฑ์ภารกิจ}{A/E1: ขั้น 2--3: บทหรือภารกิจล็อก ให้แสดงเงื่อนไขและไม่เริ่ม\newline A/E2: ขั้น 4: ภารกิจเดิมดำเนินอยู่ให้เปิดรายการเดิม; สำเร็จแล้วให้แสดงผลเดิม\newline A/E3: ขั้น 5: บันทึกล้มเหลวให้ลองใหม่โดยไม่เพิ่มความคืบหน้าซ้ำ}


\UseCaseTable{UC-12 สิ้นสุดภารกิจและรับรางวัล}{นักวิ่ง}{มีภารกิจของผู้ใช้ที่กำลังดำเนินอยู่และมีข้อมูลความคืบหน้า}{สรุปสถานะสำเร็จ บางส่วน หรือไม่สำเร็จ และบันทึกรางวัลตามเกณฑ์เพียงครั้งเดียว}{1) ผู้ใช้เลือกสิ้นสุดภารกิจ หรือเซสชันที่ผูกไว้สิ้นสุด\newline 2) ระบบตรวจหลักฐานและเงื่อนไขของภารกิจ\newline 3) ระบบบันทึกผลและรายการรางวัลสำหรับภารกิจนั้น\newline 4) ระบบแสดงสรุปความคืบหน้า หน่วยกิตที่เกี่ยวข้อง และรางวัล}{A/E1: ขั้น 2: ไม่ครบเกณฑ์ ให้ไม่สำเร็จหรือรางวัลบางส่วนเฉพาะ Side Quest ที่รองรับ\newline A/E2: ขั้น 3: เคยสรุปแล้วให้แสดงผลเดิม ไม่ให้รางวัลซ้ำ; บันทึกล้มเหลวให้ลองใหม่\newline A/E3: การสิ้นสุดภารกิจเพียงอย่างเดียวไม่ทำให้เซสชันวิ่งสิ้นสุด หากยังวิ่งอยู่ให้กลับหน้าวิ่ง; UC-10 เปิดได้หลัง UC-08 สำเร็จเท่านั้น}


\UseCaseTable{UC-13 จัดการชมรม}{นักวิ่ง; หัวหน้าชมรม}{เข้าสู่ระบบ; การแก้ไขหรือยุบต้องเป็นหัวหน้าชมรมเป้าหมาย}{สร้างชมรมใหม่ หรือบันทึกการแก้ไข/ยุบตามสิทธิ์}{1) ผู้ใช้เลือกสร้างชมรมและกรอกข้อมูล\newline 2) ระบบตรวจข้อมูลและสร้างชมรมโดยกำหนดผู้สร้างเป็นหัวหน้า\newline 3) ระบบแสดงรายละเอียดชมรมและสิทธิ์ของผู้ใช้}{A/E1: ทางเลือกจากหน้ารายละเอียด: หัวหน้าแก้ข้อมูลและระบบตรวจสิทธิ์ก่อนบันทึก\newline A/E2: ทางเลือกยุบชมรม: หัวหน้ายืนยันก่อนระบบยุบและยกเลิกสมาชิกภาพที่เกี่ยวข้อง\newline A/E3: ข้อมูลไม่ถูกต้องหรือบันทึกล้มเหลว ให้แก้ไขหรือลองใหม่\newline A/E4: ไม่ใช่หัวหน้าชมรมเป้าหมาย ให้ปฏิเสธการแก้ไขและยุบ แม้เป็นหัวหน้าชมรมอื่น}


\UseCaseTable{UC-14 เข้าร่วมหรือออกจากชมรม}{นักวิ่ง}{เข้าสู่ระบบและชมรมเป้าหมายยังใช้งานอยู่}{บันทึกคำขอเข้าร่วม สมาชิกภาพ หรือการออกจากชมรมตามเงื่อนไข}{1) ค้นหาและเลือกชมรม\newline 2) ระบบแสดงรายละเอียดและเงื่อนไขเข้าร่วม\newline 3) ผู้ใช้ยืนยันเข้าร่วม\newline 4) ระบบตรวจสมาชิกภาพและบันทึกคำขอรออนุมัติ หรือเพิ่มสมาชิกหากชมรมเปิดรับทันที\newline 5) ระบบแสดงสถานะล่าสุด}{A/E1: ขั้น 4: มีสมาชิกภาพหรือคำขอเดิมให้แสดงรายการเดิม\newline A/E2: ทางเลือกออกจากชมรม: สมาชิกยืนยันแล้วระบบยุติสมาชิกภาพ; หัวหน้าต้องโอนบทบาทหรือยุบชมรมก่อนออก\newline A/E3: ชมรมปิดรับ ถูกยุบ หรือบันทึกล้มเหลว ให้แจ้งเหตุผลและไม่แสดงว่าเข้าร่วมสำเร็จ\newline A/E4: นโยบายเปิดรับทันทีหรืออนุมัติต้องกำหนดก่อนใช้งานจริง; flow นี้เป็นข้อกำหนด ไม่ใช่ผลยืนยันการทดสอบสองบัญชี}


\UseCaseTable{UC-15 จัดการสมาชิกในชมรม}{หัวหน้าชมรม}{เข้าสู่ระบบและเป็นหัวหน้าของชมรมเป้าหมาย}{สมาชิกภาพหรือบทบาทถูกปรับตามคำสั่งที่ได้รับอนุญาต}{1) เปิดรายชื่อสมาชิกหรือคำขอของชมรมตนเอง\newline 2) เลือกอนุมัติ/ปฏิเสธคำขอ นำสมาชิกออก หรือโอนหัวหน้าให้สมาชิกที่มีสิทธิ์\newline 3) ระบบขอคำยืนยันและตรวจสิทธิ์ล่าสุด\newline 4) ระบบบันทึกการเปลี่ยนแปลงและแสดงผล}{A/E1: ขั้น 3: ยกเลิกให้คงข้อมูลเดิม; ไม่มีสิทธิ์ให้ปฏิเสธ\newline A/E2: ขั้น 4: สมาชิกหรือคำขอเปลี่ยนไปแล้ว ให้โหลดใหม่\newline A/E3: การโอนหัวหน้าต้องสำเร็จทั้งการแต่งตั้งคนใหม่และเปลี่ยนบทบาทคนเดิม ไม่ปล่อยชมรมไม่มีหัวหน้า\newline A/E4: การจัดการสมาชิกระหว่างหลายบัญชียังไม่มีหลักฐานทดสอบครบในรายงาน}


\UseCaseTable{UC-16 จัดการรายชื่อเพื่อน}{นักวิ่ง}{เข้าสู่ระบบ}{คำขอหรือความสัมพันธ์เพื่อนเปลี่ยนตามการกระทำที่ยืนยัน}{1) ค้นหาผู้ใช้และเลือกส่งคำขอเป็นเพื่อน\newline 2) ระบบตรวจว่าไม่ใช่ตนเองและไม่มีคำขอหรือความสัมพันธ์ซ้ำ\newline 3) ระบบบันทึกคำขอและแสดงสถานะรอตอบรับ}{A/E1: ทางเลือกตอบรับ/ปฏิเสธ: ผู้รับคำขอเลือกและระบบตรวจสิทธิ์ก่อนเปลี่ยนสถานะ\newline A/E2: ทางเลือกยกเลิกคำขอหรือลบเพื่อน: ผู้มีสิทธิ์ยืนยันก่อนบันทึก\newline A/E3: ขั้น 2: ขอเป็นเพื่อนตนเองหรือซ้ำ ให้ปฏิเสธและแสดงสถานะเดิม\newline A/E4: บันทึกล้มเหลวให้ลองใหม่; การตอบรับระหว่างสองบัญชียังต้องทดสอบเพิ่มเติม}


\UseCaseTable{UC-17 ซื้อไอเทม}{นักวิ่ง}{เข้าสู่ระบบและสินค้ายังเปิดขาย}{หักเหรียญและเพิ่มไอเทมในคลังสำเร็จเป็นรายการซื้อเดียวกัน}{1) เปิดร้านและเลือกสินค้า\newline 2) ระบบแสดงราคาและยอดเหรียญ\newline 3) ผู้ใช้ยืนยันซื้อ\newline 4) ระบบตรวจยอดเหรียญ สิทธิ์ซื้อ และรายการที่เคยซื้อ\newline 5) ระบบหักเหรียญ เพิ่มไอเทม และบันทึกประวัติการซื้อ\newline 6) ระบบแสดงผลและยอดคงเหลือ}{A/E1: ขั้น 3: ยกเลิกไม่เปลี่ยนยอด\newline A/E2: ขั้น 4: เหรียญไม่พอหรือมีไอเทมที่ซื้อซ้ำไม่ได้ ให้ปฏิเสธ\newline A/E3: ขั้น 5: บันทึกล้มเหลวต้องไม่เหลือการหักเหรียญโดยไม่มีไอเทม; ส่งคำขอเดิมซ้ำไม่หักเหรียญซ้ำ}


\UseCaseTable{UC-18 จัดการอวตาร์}{นักวิ่ง}{เข้าสู่ระบบและมีอวตาร์ของตนเอง}{บันทึกการแต่งตัวด้วยไอเทมที่ผู้ใช้ครอบครอง}{1) เปิดอวตาร์และคลังไอเทม\newline 2) ระบบแสดงไอเทมที่ใช้ได้กับแต่ละช่อง\newline 3) ผู้ใช้เลือกใส่หรือถอดไอเทมและยืนยัน\newline 4) ระบบตรวจความเป็นเจ้าของและความเข้ากันได้ของช่อง แล้วบันทึก\newline 5) ระบบแสดงอวตาร์ล่าสุด}{A/E1: ขั้น 3: ยกเลิกให้คงการแต่งตัวเดิม\newline A/E2: ขั้น 4: ไม่มีไอเทมหรือช่องไม่ตรง ให้ปฏิเสธ\newline A/E3: บันทึกล้มเหลวให้ลองใหม่; ผลทดสอบที่มีครอบคลุมการอ่านคลัง ยังไม่ยืนยันการแก้ไขครบทุกกรณี}


\UseCaseTable{UC-19 ดูตารางจัดอันดับ}{นักวิ่ง}{เข้าสู่ระบบและมีสิทธิ์ดูขอบเขตอันดับที่เลือก}{แสดงอันดับตามตัวชี้วัด ช่วงเวลา และกลุ่มที่ระบุ}{1) เปิดตารางอันดับและเลือกกลุ่มเพื่อนหรือชมรมที่เข้าถึงได้\newline 2) ระบบแสดงตัวชี้วัด ช่วงเวลา และเกณฑ์กรณีคะแนนเท่ากัน\newline 3) ระบบคำนวณหรืออ่านอันดับที่ตรงเงื่อนไข\newline 4) ระบบแสดงอันดับและข้อมูลสรุปที่อนุญาตให้เผยแพร่}{A/E1: ขั้น 1: ไม่มีสิทธิ์ดูชมรม ให้ปฏิเสธ\newline A/E2: ขั้น 3: ไม่มีข้อมูล ให้แสดงสถานะว่าง; โหลดล้มเหลวให้ลองใหม่\newline A/E3: ตัวชี้วัด ช่วงเวลา และเกณฑ์คะแนนเท่ากันต้องกำหนดก่อนพัฒนา และยังไม่มีผลทดสอบ Leaderboard ในรายงาน}


\UseCaseTable{UC-20 เชื่อมต่อและนำเข้าการวิ่งจาก Strava}{นักวิ่ง; Strava}{เข้าสู่ระบบและมีบัญชี Strava; เป็นขอบเขตการออกแบบสำหรับการต่อยอด}{นำเข้ากิจกรรมที่ได้รับอนุญาตโดยไม่สร้างกิจกรรมซ้ำ}{1) ผู้ใช้เลือกเชื่อมต่อ Strava\newline 2) Strava ให้ผู้ใช้อนุญาตสิทธิ์และระบบตรวจผลยืนยัน\newline 3) ผู้ใช้สั่งนำเข้าและระบบขอกิจกรรมในขอบเขตที่อนุญาต\newline 4) ระบบตรวจรหัสกิจกรรมจากแหล่งข้อมูลและบันทึกเฉพาะรายการที่ยังไม่มี\newline 5) ระบบแสดงผลนำเข้าและแหล่งที่มาของข้อมูล}{A/E1: ขั้น 2: ยกเลิกหรือไม่อนุญาต ให้คงสถานะไม่เชื่อมต่อ\newline A/E2: ขั้น 3: สิทธิ์หมดอายุ ถูกถอน หรือบริการไม่พร้อม ให้เชื่อมต่อใหม่หรือลองภายหลัง\newline A/E3: ขั้น 4: รายการเดิมให้ข้าม ไม่เพิ่มระยะทางหรือรางวัลซ้ำ\newline A/E4: ทางเลือกยกเลิกเชื่อมต่อ: ยุติการนำเข้าครั้งถัดไป; การลบประวัติต้องเป็นคำสั่งแยก\newline A/E5: ยังไม่ใช้ข้อมูลนำเข้าเพื่อให้หน่วยกิตหรือรางวัลจนกำหนดเกณฑ์ตรวจสอบครบ}

